\documentclass[a4paper,12pt]{article}

% 基础宏包配置
\usepackage{ctex}      % 中文支持
\usepackage{geometry}  % 页面布局
\usepackage{titlesec}  % 标题格式
\usepackage{fancyhdr}  % 页眉页脚
\usepackage{xcolor}    % 颜色支持
\usepackage{setspace}  % 行间距

% 页面边距设置
\geometry{left=2.5cm, right=2.5cm, top=3cm, bottom=3cm}
\onehalfspacing % 设置1.5倍行距，提升阅读体验

% 颜色定义
\definecolor{primaryBlue}{RGB}{0, 56, 107}
\definecolor{darkGrey}{RGB}{60, 60, 60}

% 标题格式定制
\titleformat{\section}
{\normalfont\Large\bfseries\color{primaryBlue}}{\thesection}{1em}{}
\titleformat{\subsection}
{\normalfont\large\bfseries\color{darkGrey}}{\thesubsection}{1em}{}

% 页眉页脚
\pagestyle{fancy}
\fancyhf{}
\lhead{\small \textbf{2024年度ESG报告}}
\rhead{\small 可持续发展与社会责任}
\cfoot{\thepage}

\title{\textbf{\huge 2024年度环境、社会及治理（ESG）报告}\\ \vspace{0.5em} \large 科技向善，绿动未来}
\author{本公司董事会办公室}
\date{2024年3月}

\begin{document}

\maketitle

\section{战略引领：新质生产力与稳健治理}

2024年，在全球汽车产业加速变革与国家“双碳”战略深入推进的背景下，本公司始终坚持“三直四化”——即电动化、智能网联化、高端化、国际化的发展战略。我们深刻理解，所谓“新质生产力”，不仅意味着技术的革新，更代表着一种绿色、高效、可持续的增长模式。这一年，我们积极拥抱变化，将ESG理念融入企业经营的毛细血管，致力于成为全球领先的客车及出行解决方案提供商。

\subsection{稳健的经营绩效与股东回报}

在复杂多变的宏观经济环境下，公司依托技术积累与全球化布局，交出了一份稳健的答卷。2024年度，集团实现营业收入372.18亿元，同比增长37.63\%；归属于母公司股东的净利润达到41.16亿元。我们深知，企业的长远发展离不开股东的信任与支持。为了践行长期主义，让投资者共享企业发展红利，公司在报告期内极其罕见地实施了两次现金分红，全年累计发放现金股利达44.28亿元，分红总额甚至超过了当年的净利润，充分体现了我们回报股东的诚意与实力。

\subsection{合规底线与供应链生态}

合规经营是企业立足的基石。我们引以为傲的是，公司已连续第13年获得交易所信息披露工作的“A级”评价，这在资本市场中是一项极高的荣誉。同时，我们顺利通过了ISO 37301合规管理体系认证，标志着公司的内部管理流程已完全接轨国际标准。

在供应链管理方面，我们致力于带动产业链上下游共同发展。目前，我们与530家供应商保持着紧密的合作关系，其中20\%的合作伙伴来自河南本地，有力拉动了区域装备制造业的集群发展。我们不仅关注产品的交付，更关注商业道德。截至2024年底，我们已与包括潜在供应商在内的超过900家合作伙伴签署了廉洁协议，通过严格的制度约束，共同营造风清气正、阳光透明的商业生态。

\section{社会责任：创新驱动与人文关怀}

我们坚信，科技创新的最终目的是为了服务于人。无论是通过技术手段提升车辆的安全与能效，还是通过公益活动回馈社区，我们始终将“人”放在核心位置。

\subsection{研发创新：用科技定义美好出行}

在“软件定义汽车”的时代，客车早已不再是简单的交通工具，而是集成了高精尖技术的移动终端。2024年，我们在研发领域投入了17.88亿元，占营业收入的4.80\%。这笔巨大的投入支持着我们一支由3,505名研发人员组成的庞大团队，其中包括16位博士领军人才。辛勤的耕耘换来了丰硕的成果，全年我们新增授权专利176项，其中发明专利占比过半，达到114项；同时，我们还累计参与制定了317项国家及行业标准，始终走在行业技术规范的前沿。

为了让公众更好地理解我们的技术突破，我们将几项晦涩的工程技术进行了通俗化的解读。首先是自主研发的“YOS车机操作系统”。这就像智能手机的安卓系统一样，它实现了汽车软件与硬件的“解耦”。过去，车辆功能出厂即定型，而现在，通过OTA（空中下载技术），用户可以像升级手机APP一样远程升级车辆系统，让车常用常新。

其次是颠覆性的“C架构”。传统的商用车内部有几十个控制器，线束错综复杂。我们的C架构通过“中央计算+区域控制”的模式，相当于给汽车装上了一个超级大脑和几个执行组长，不仅大幅减少了线束的使用和车辆自重，更让车辆的反应速度和智能化水平有了质的飞跃。

在核心零部件方面，我们研发的“多合一动力域控制器”应用了先进的碳化硅（SiC）模块。相比传统的硅基芯片，碳化硅的开关损耗降低了50\%至70\%，这直接提升了能量转化效率，让每一度电都能驱动车辆跑得更远。针对电动车冬季续航缩水的行业痛点，我们推出了“多源超低温热泵”技术。它能巧妙地从寒冷的空气甚至电机余热中“榨取”热量，实测显示，在低温环境下其采暖节能效果高达30\%，有效缓解了用户的冬季里程焦虑。

此外，我们在智能网联领域也取得了里程碑式的进展。公司的L4级自动驾驶巴士已累计安全运营超过1,700万公里，并成功获批成为国家首批智能网联汽车准入试点企业，这标志着无人驾驶技术正逐步从实验室走向大规模商业化应用。

\subsection{员工权益与安全守护}

员工是公司最宝贵的财富。截至2024年底，集团共有员工16,707人，其中女性员工占比为10.02\%。我们致力于为每一位员工提供安全、健康的工作环境。在过去的一年里，我们在安全生产上的投入毫不吝啬：投入5,500余万元用于采购高标准的劳动防护用品，投入2,200余万元用于车间除尘、降噪和通风设施的改善，投入180余万元用于全员安全培训。

这一系列“真金白银”的投入换来了显著的成效。全年我们对7.9万人次进行了安全教育培训，实现了无新增职业病的目标，切实守护了每一位劳动者的身心健康。

\subsection{社区贡献：传递温暖与绿色}

我们不仅造好车，更努力做一个好邻居。2024年，公司对外捐赠及公益项目总投入超过4,000万元，重点聚焦于交通安全教育和生态环境保护。

在儿童安全领域，我们的品牌公益项目“儿童交通安全公益行”这一年走过了全国13个省市。在华东某市的小学操场上，我们将一辆真实的重型车辆开了进去，并在车辆周围设置了模拟障碍物。我们邀请孩子们坐进驾驶室，利用后视镜观察四周。当孩子们亲眼发现某些区域在镜子里完全看不见时，他们对“盲区”有了最直观、最震撼的认识。这种体验式的教育，比枯燥的说教更能让近万名受益儿童记住如何避开危险。

作为一家国际化企业，我们在出海的同时也不忘回馈全球社区。在南美某国，针对当地干旱化的环境问题，我们发起了“零碳森林”计划，种植了大量适宜当地气候的树种。截至目前，我们在全球范围内已累计植树36,000余棵。这片不断生长的森林，不仅能持续吸收二氧化碳，更改善了当地的生态环境，成为了中国企业在海外履行社会责任的一张绿色名片。

\section{环境绩效：打造绿色制造标杆}

面对全球气候变化的挑战，我们承诺在生产运营的每一个环节都践行低碳理念。通过精细化的能源管理和技术改造，我们正在将“绿色制造”变为现实。

\subsection{碳排放与能源管控}

2024年，在公司营收大幅增长的同时，我们通过技术手段有效控制了碳排放强度。数据显示，全年范围1（直接能源燃烧）排放量为57,330.53吨二氧化碳，范围2（外购电力）排放量为122,786.63吨二氧化碳。通过计算，我们的碳排放强度控制在0.07吨二氧化碳/万元营收的水平。综合能源消费量为34,560.94吨标准煤。

为了降低这些数字，我们实施了扎实的节能技改行动。全年我们落地了7项大型节能改造项目，通过升级高能效设备和优化工艺流程，不仅提升了生产效率，更实现了年节电1,258兆瓦时、年节气11.34万立方米的实质性减排。此外，我们高度重视余热资源的回收利用，全年共回收利用热能约29,306吉焦，这些原本会被浪费的热量被重新用于厂区供暖和工业加热，大幅减少了化石能源的消耗。

\subsection{清洁能源与资源循环}

我们正在积极调整能源结构，让清洁能源成为生产动力的重要组成部分。在新能源智造基地，屋顶光伏电站全年发电量达到155.68万千瓦时，这些自产的“绿电”直接减少了约835吨碳排放。同时，我们还通过电力市场采购了13,580兆瓦时的绿色电力，进一步提升了可再生能源的使用比例。

在水资源管理方面，我们把水看作可循环流动的资产。2024年，公司新取水量为141.25万吨，但通过完善的中水回用系统，我们的重复用水量高达6,465.36万吨，重复用水率达到惊人的97.8\%。这意味着，绝大部分工业用水都在工厂内部实现了闭环循环。

在废弃物治理上，我们也坚持“变废为宝”和合规处置并重。全年产生的一般工业固废中，有61,485.76吨被回收利用，让废钢、废纸等资源重新回到了生产循环。对于产生的4,110.84吨危险废物，我们严格委托有资质的单位进行了100\%的合规处置，确保不留任何环境隐患。此外，我们在物流环节大力推行循环包装，目前国产零部件的周转包装占比已达到89\%，这种可重复使用的塑料箱大幅替代了传统的一次性纸箱，有效减少了物流垃圾的产生。

\section{未来展望}

2024年是本公司深化ESG实践的重要一年，但绿色转型是一场没有终点的马拉松。展望未来，我们将继续以“三直四化”战略为指引，在氢能利用、电池回收、供应链碳管理等前沿领域持续探索。我们相信，通过不懈的技术创新和责任担当，我们一定能为全人类创造更加绿色、智能、美好的出行未来。

\vspace{2cm}

\end{document}